\section{Step 28 Compare one update with several updates}\label{sha-s28}

Supply the same 65 bytes in one call and in parts of fifty-five and ten. Compare the state, total, used count and used buffer prefix before finalization. These are the meaningful fields. A correct implementation may leave different unused buffer or schedule bytes after the two paths.

\subsection{Observe and compare}
Use recipes \texttt{s28-1}, \texttt{s28-2}, \texttt{s28-3}. Prepare each in a separate source image from the current cumulative program. An internal checkpoint is inspected before its temporary or normal return; otherwise pause at \texttt{main\_done}.

Use the completed entries available at this point; keep their inputs and observation boundary explicit. The surrounding explanation identifies the state being compared.

\subsection{Expected observation}
Both paths represent sixty-five bytes, with one buffered byte and the same state after the first full block.

\Needspace{9\baselineskip}
\subsection{Check your understanding}
Which context bytes must represent the same tail?

\begin{itemize}
  \item \textbf{A} The entire unused schedule.
  \item \textbf{B} Only the used buffer prefix.
  \item \textbf{C} Every byte ever used by any earlier call.
\end{itemize}
